\documentclass[a4paper,12pt]{article}
	\usepackage[ansinew]{inputenc}
	\usepackage{ae}
	\usepackage[brazil]{babel}
	\usepackage{lipsum}
	\usepackage{xcolor}
	\usepackage{graphicx}	
	\usepackage{subfigure} % O pacote sugfigure facilita a criação de estruturas como fig.1, fig.1(a) etc
	
	\graphicspath{
		{Figuras/EPS/}
		{Figuras/JPG/}
		{Figuras/PDF/}
		{Figuras/PNG/}
		}
	
	\newcommand{\comando}[1]{{\normalfont\textbackslash\color{blue!50!black}#1}}
\newcommand{\pacote}[1]{{\normalfont\sffamily\color{orange!50!black}#1}}
\newcommand{\ambiente}[1]{{\normalfont\sffamily\color{green!50!black}#1}}
\newcommand{\contador}[1]{{\normalfont\sffamily\color{red!50!black}#1}}
\newcommand{\foreign}[1]{\textsl{#1}}
\newcommand\linhabase{\noindent\makebox[0pt][l]{\color{red}\rule{\textwidth}{0.1pt}}}
	
\begin{document}

	\lipsum[1]
	
	\begin{figure}[tb]
		\hfill\subfigure
				[Esta é a legenda da figura à esquerda, figura~\ref{fig:figA}. Observe que com o pacote \pacote{subfigure} não é mais necessário utilizar \ambiente{minipage}s, tornando a sintaxe mais simples.]
				{\label{fig:figA}\includegraphics[width=0.45\textwidth]{Tei-Gi}}% <-- Para quê serve este comentário?
		\hfill
		\subfigure
				[Esta é a legenda da figura à direita, figura~\ref{fig:figB}. O tamanho da fonte, o alinhamento do texto e outras configurações das legendas são parametrizadas pelo pacote \pacote{subfigure}, de modo que podem ser facilmente modificadas.]
				{\label{fig:figB}\includegraphics[width=0.45\textwidth]{Tei-Gi}}% <-- Para quê serve este comentário?
		\hfill\null % <-- Você pode trocar \hfill\null por \hspace*{\fill}
		
		\caption{esta é a legenda comum às duas figuras. Observe que se atribuirmos um rótulo (comando \comando{label}) a cada uma das sub-figuras (no comando \comando{subfigure}) e à figura como um todo (no ambiente \ambiente{figure}), podemos utilizar o comando \comando{ref} para gerar as referências \ref{fig:figA}, \ref{fig:figB} e~\ref{fig:figAeB}.}
		\label{fig:figAeB}
	\end{figure}
		
	% Tarefas extras:
	% 1) Modifique os \hfill acima de modo que as duas figuras distem entre si de 2em (dica: use o comando \qquad).
	% 2) Faça outra alteração de modo que a figura à esquerda fique encostada na margem esquerda e a da direita, na margem direita.
	% 3) Na sub-figura deste exemplo, a legenda aparece assim:
	%
	%           (a) legenda legenda legenda legenda
	%           legenda  legenda   legenda  legenda
	%           legenda  legenda  legenda  legenda.
	%
	% Consulte a documentação do pacote subfigure e altere este arquivo de modo que a legenda apareça assim:
	%
	%           (a) legenda legenda  legenda legenda
	%               legenda  legenda legenda legenda
	%               legenda legenda legenda legenda.
	%
	% Isto é, faça com que (a) e (b) fique "fora" do corpo da legenda. Dica: veja a figura 8 e a seção 2 da documentação.
	
\end{document}